\section{Policy improvement, stopping and the learning graph}\label{sec:policy119}
The certificate separates three choices: which continuation to acquire, which
feasible action to propose, and whether to spend more work on that action.
The first two determine a candidate complete policy. A tighter enclosure of
that same policy changes the information available for the third choice.
Actual improvement requires a comparison of original policy values.

\paragraph{Completed policies and foregone gain.}
Let a completed feasible incumbent have a verified cost bracket
$[L_B,U_B]$ and let $O_B$ be a valid lower bound on the original optimum.
A completed trial with upper cost $U_T<L_B$ delivers a strict original-policy
improvement, provided it also meets the unchanged qualification targets.
The improvement is at least $L_B-U_T$. The remaining possible gain from any
further feasible policy is at most
\[
 G_B=\max\{0,U_B-O_B\}.
\]
Thus stopping when $G_B\le\eta$ has an explicit original-objective opportunity
cost. For a monetary value $v$ per unit of original loss and a lower bound $c$
on an additional computation's cost, $c\ge vG_B$ rules out a strictly positive
net gain from that computation. The converse inequality does not guarantee
that a useful improvement will be found. The complete original theorems,
interface sharpness and recursive risk-distorted composition are retained in
Development Sections~\ref{sec:improvement109} and~\ref{sec:budget110}, with
full proofs in the supplement. Their R110 experiments retain their own actual
policy gains and positive foregone gains; they are not new R119 observations.

\paragraph{An action comparison can share a risk witness.}
For bounded original continuations $X,Y$ and one common feasible dual law $q$,
write $k_X(q)=\rho(X)-\{\E_qX-\psi(q)\}\ge0$. Then
\[
 \rho(Y)-\rho(X)=\E_q(Y-X)+k_Y(q)-k_X(q).
\]
Supplement Theorem~\ref{thm:coupled119} proves a computable action contrast
using the two original-risk defects, with an atom-safe quantile-band remainder.
Its implemented-future allowance is charged once on each side of the contrast.
Proposition~\ref{prop:nearactive119} proves that a continuation contrast of
oscillation $w>0$ only needs dual laws whose reference optimality slack is at
most $w$, without assuming the supremum is attained. This preserves the signed
recursive policy difference instead of discarding the reference slack.
The integral of a difference of convex continuations still needs a valid
separate enclosure. This mathematical interface advances the independently
frozen R111--R112 question; no current numerical result is attributed to a
paired algorithm that has not been implemented.

\paragraph{The actor and critic are a specified computational graph.}
For a smooth critic $u_\theta$ and a differentiable feasible actor $a_\omega$,
the Bellman residual includes the derivative through the actor and through
every parameter-dependent continuation and transition. Differentiating the
half-squared residual gives $\E[r_{\theta,\omega}\nabla r_{\theta,\omega}]$
under the integrability and differentiation hypotheses in Supplement
Section~\ref{supp:foundations119}. Holding a target fixed defines a different
update and requires its own convergence argument. Discrete original-law
qualification treats a realized rounded action as a feasible decision and
does not differentiate its numerical rounding operation.

For the Brownian Markov recursive-utility interface, the critic graph is
$z_u=\sigma^{\mathsf T}\nabla u$. The driver, admissible controls, terminal
condition and stopping boundary belong to that same recursive criterion.
Supplement Proposition~\ref{prop:stoppedbsde119} proves the stopped-domain
residual account, including its stability factor and boundary discrepancy.
A general ReLU critic is not a classical $C^{1,2}$ It\^o test function merely
because its ordinary Hessian vanishes away from kinks. The discrete NBO
construction remains valid under its own hypotheses. Fully coupled NDU/FBSDE
systems and other utility objectives require their own well-posedness and
verification premises; they do not inherit the investment experiment's
numerical ranking by a change of terminology.
